% Generated draft: requires mathematical and source review.
\documentclass[11pt]{article}
\usepackage[margin=25mm]{geometry}
\usepackage{fontspec}
\setmainfont{Latin Modern Roman}
\usepackage{amsmath,amssymb,mathtools,unicode-math}
\setmathfont{Latin Modern Math}
\usepackage{enumitem,needspace,xurl,hyperref}
\setlength{\parindent}{0pt}
\setlength{\parskip}{5pt}
\setlength{\emergencystretch}{3em}
\newcommand{\N}{\mathbb{N}}
\newcommand{\Z}{\mathbb{Z}}
\newcommand{\Q}{\mathbb{Q}}
\newcommand{\R}{\mathbb{R}}
\newcommand{\C}{\mathbb{C}}
\newcommand{\F}{\mathbb{F}}
\newcommand{\A}{\mathbb{A}}
\newcommand{\PP}{\mathbb{P}}
\newcommand{\E}{\mathbb{E}}
\newcommand{\eps}{\varepsilon}
\newcommand{\dd}{\,\mathrm{d}}
\begin{document}
\section*{UnsolvedMath ID: 30004720}
\textbf{Source problem code:} catalogue entry
\subsection{Definitions and mathematical statement}
Let \(g\) be a continuous Riemannian metric, meaning a continuous field of positive-definite symmetric bilinear forms. The source asks about two notions of a lower scalar-curvature bound for such metrics: the Ricci-flow-based \(\beta\)-weak notion and Gromov's polyhedral formulation. Their definitions, the admissible range of \(\beta\), and the hypotheses on the underlying manifold are not supplied. Consequently, the following questions retain these named notions as unresolved source-dependent terms.

\begin{enumerate}
\item For a scalar lower bound \(k\in\R\), determine whether the Ricci-flow-based \(\beta\)-weak condition and Gromov's polyhedral condition are equivalent. The source requests both local and global comparisons, but does not specify the precise local condition, the global hypotheses, or the quantification over \(\beta\).
\item Consider a continuous Riemannian metric near \(0\in\R^n\). For \(x\ne0\), put \(r=|x|\) and \(\theta=x/r\in S^{n-1}\), where \(|\cdot|\) is the Euclidean norm. Suppose its coordinate components have the stated asymptotic form
\[
 g_{ij}(x)=\delta_{ij}+r^2G_{ij}(\theta)+O\!\left(|x|^{2+\eta}\right)
 \qquad (x\to0),
\]
where \(\delta_{ij}\) are the Euclidean metric components and the coefficient functions \(G_{ij}\) satisfy
\[
 \sum_{i,j=1}^{n}\theta^i\theta^jG_{ij}(\theta)=0
 \qquad\text{for every }\theta\in S^{n-1}.
\]
Give an explicit necessary-and-sufficient characterization of when \(g\) has nonnegative scalar curvature at \(0\) in the Ricci-flow-based \(\beta\)-weak sense.
\end{enumerate}

The second question concerns the stated weak condition at a point, not an asserted existence of classical scalar curvature there. The source does not specify the allowed dimensions, the value or range of \(\eta\), the regularity of \(G_{ij}\), or the precise control required of the remainder. In particular, positivity of \(\eta\) is not explicitly supplied. The notation \(r=|x|\) is the radial interpretation of the displayed source expression. A fully quantified formulation requires these missing hypotheses and the exact definitions of both curvature notions; no equivalence or characterization is asserted here.
\subsection{Short English statement}
The problem asks whether two proposed lower scalar-curvature conditions for continuous metrics agree, both locally and globally. It also asks for an explicit criterion for weak nonnegative scalar curvature at the origin for metrics with the displayed quadratic angular expansion. The definitions of the curvature conditions and several parameter and regularity hypotheses are missing, so this remains a qualified statement-only draft.
\subsection{Sources}
\begin{itemize}
\item Supplied local UnsolvedMath registry record, permanent ID \(30004720\), source problem code \texttt{catalogue entry}, titled \emph{Equivalence of Weak Scalar-Curvature Notions for Continuous Metrics}, in the \emph{Oberwolfach Reports Open Problems} collection. No individual proposer is supplied; the proposed year is recorded as \(2021\).
\item Supplied report citation: \emph{Analysis, Geometry and Topology of Positive Scalar Curvature Metrics} (2022), Oberwolfach Reports, DOI \texttt{10.4171/owr/2021/30}, pp.\ 1579--1649. \url{https://doi.org/10.4171/owr/2021/30}. The registry reports that Questions 8 and 10 were combined and that the full asymptotic metric expression and tangentiality condition were restored. Its label \texttt{corrected\_verified} and its claim that the report was inspected are supplied curation, not independent verification in this draft.
\item Supplied literature reference: Burkhardt-Guim, ``pointwise scalar bounds.'' \url{https://arxiv.org/abs/1907.13116}.
\item Supplied literature reference: Gromov, ``scalar curvature lectures.'' \url{https://arxiv.org/abs/1908.10612}.
\item Literature assessment dated exactly \texttt{2026-08-22}: the supplied curation labels the problem open and states that the Ricci-flow and approximation formulations are identified, while equivalence with Gromov's polyhedral formulation and the requested explicit local characterization remain open. This is supplied dated curation derived from the OWR record, not an independently verified literature assessment. Time-sensitive status and cited claims require independent verification.
\end{itemize}
\end{document}
